\scanpage{134}
\mspage{-III- \handwritten{3} -}
\begin{flushright}number \handwritten{383}\end{flushright}
for $G^{\unsure{o|0|\circ}}$ is the plane $P$ spanned by $e_i$ and $e_j$. The hyperplanes of the form $g(H_i)$ and $g(H_j)$, for $g\in G^{\unsure{o|0|\circ}}$, all contain the subspace $P'$ orthogonal to $P$, and form a subset $\underline{H}^{\unsure{o|0|\circ}}$ of $\underline{H}$ which has, with respect to $G^{\unsure{o|0|\circ}}$, properties 1) and 2) above. Since $\underline{H}^{\unsure{o|0|\circ}}\handwritten{\subset}\underline{H}$, there exists a chamber $C^{\unsure{o|0|\circ}}$ defined by $\underline{H}^{\unsure{o|0|\circ}}$ and containing $C$ ; it is clear that $C^{\unsure{o|0|\circ}}$ is one of the dihedra having $H_i$ and $H_j$ as walls. Thus $C^{\unsure{o|0|\circ}}\handwritten{\cap}P$ is an angular sector, bounded by two half-lines $D_i\handwritten{\subset}H_i$ and $D_j\handwritten{\subset}H_j$. Setting $\handwritten{\vartheta}=(\handwritten{\widehat}{D_i,D_j})$, $s_j s_i$ is the rotation through $2\handwritten{\vartheta}$ about $P'$, and the transforms of $C^{\unsure{o|0|\circ}}$ under $G^{\unsure{o|0|\circ}}$ are its transforms under rotations through angle $n\handwritten{\vartheta}$ $(n\in Z)$ ; since two distinct transforms of $C^{\unsure{o|0|\circ}}$ are disjoint, it follows from \S\ 1, number 3, lemma 9 that $\handwritten{\vartheta}$ is of the form $\handwritten{\pi}/m$ $(m\in Z)$ ; we necessarily have $m=m_{ij}$. By definition of $e_i$ and $e_j$, and with a suitable orientation of $P$, we have $(\handwritten{\widehat}{D_i,e_i})=-\handwritten{\pi}/2$, and $(\handwritten{\widehat}{D_j,e_j})=+\handwritten{\pi}/2$. Hence, by Chasles's formula (Algebra IX, \S\ 10, number 3, (9)), $(\handwritten{\widehat}{e_i,e_j})=(\handwritten{\widehat}{e_i,D_i})+(\handwritten{\widehat}{D_i,D_j})+(\handwritten{\widehat}{D_j,e_j})=\handwritten{\pi}+(\handwritten{\widehat}{D_i,D_j})$. We deduce $(e_i\handwritten{|}e_j)=\cos(\handwritten{\widehat}{e_i,e_j})=-\cos(\handwritten{\widehat}{D_i,D_j})=-\cos(\handwritten{\pi}/m_{ij})$. The second assertion follows from the monotonicity of the cosine function on $[0,\handwritten{\pi}/2]$ and the formula $\cos(\handwritten{\pi}/3)=1/2$ (reference impossible to find !!).

PROPOSITION 3.- \uline{The vectors} $e_1,\ldots,e_m$ \uline{form a basis of the essential subspace} $W$ \uline{of} $E$.

Indeed the set $C'$ consisting of $0$ and of the $y\in E$ such that $(y\handwritten{|}x)>0$ for every $x\in C$ is a pointed convex cone, and we have $C'\handwritten{\cap}(-C')=\handwritten{\{}0\handwritten{\}}$ ; thus $C'$ is the set of positive elements for an order relation making $E$ an ordered vector space (Topological Vector Spaces, II, \S\ 1, number 5). By (2) we have $e_i>0$ for every $i$, and by proposition 2, we have $(e_i\handwritten{|}e_j)\handwritten{\leq}0$ for $i\handwritten{\ne}j$. Lemma 6 of \S\ 1, number 2 then shows that the $e_i$ are linearly independent. On the other hand it is clear that the $e_i$ belong to $W$. Finally if an element of $E$ is orthogonal to $e_1,\ldots,e_m$, it is fixed by the
% archive sheet number: 133
